\begin{samepage}
\bigskip
\pdfbookmark[0]{Magnetic Observations. 1835}{magnetische_beobachtungen_1835}
\begin{center}
\hrule\vspace{2mm}
Göttingische gelehrte Anzeigen. 7 March 1835.
\vspace{2mm}\hrule\vspace{5mm}
\end{center}
At the meeting of the Royal Society on 14 February, Court Councillor \textsc{Gauss} reported on the observations made at the magnetic observatory and on related observations made elsewhere. We give an extract from that report here, which may be taken as
\begin{center}\itshape a continuation of the notice given on 9 August 1834.\end{center}
\end{samepage}

\Needspace{13\baselineskip}
\SourceLeaf{535}The twice-daily records of the position of the needle have continued without interruption and now cover nearly a full year. The monthly means since July of last year were:
\begin{center}
\begin{tabular}{rl|rrr|rrr}
 & & \multicolumn{3}{c|}{8 a.m.} & \multicolumn{3}{c}{1 p.m.}\\\hline
1834 & August & $18^0$ & $38'$ & $48''1$ & $18^0$ & $49'$ & $11''0$\\
 & September & & $36$ & $58.4$ & & $46$ & $32.3$\\
 & October & & $37$ & $18.4$ & & $44$ & $47.2$\\
 & November & & $37$ & $38.4$ & & $43$ & $4.3$\\
 & December & & $37$ & $54.8$ & & $41$ & $32.7$\\
1835 & January & & $37$ & $51.5$ & & $42$ & $14.4$
\end{tabular}
\end{center}

The observations agreed upon for specified days at short, uninterrupted intervals have been carried out on four principal occasions since the last notice, apart from a few special additional occasions. In recent months, some changes to the arrangements for these observations have been made and announced elsewhere. Participation has already spread and will soon spread further. Some very remarkable results have already come from these observations, similar to those mentioned in the earlier report. A graph showing the agreement among the observations of 1 and 2 October and of 29 and 30 November in Göttingen, Leipzig and Berlin will appear shortly in \textsc{Poggendorff}'s Annalen der Physik. Even more remarkable, however, is the agreement between the observations of 5 and 6 November in Copenhagen and Milan in all the numerous and strikingly large fluctuations. A drawing of these will also be given elsewhere. Here we enter a world of mysterious forces of nature, whose wonderfully changing play reveals itself in the same way, at the same instant and down to the smallest nuances, across half the breadth of Europe. This baseline still seems far too short to measure the range of their action.

The facilities here for magnetic observations have meanwhile been substantially extended in several ways. In some observations, when great precision is required, it is essential to use a second apparatus at some distance from the main one, so that hourly changes in the magnetic force can be taken into account. For this purpose, a large magnetic bar has hung in the astronomical observatory since August of last year, after the instruments used in 1832 \SourceLeaf{536}were handed over to the physics cabinet. Its other fittings are very much like those of the bar in the magnetic observatory. The magnetic bar in the astronomical observatory is likewise made of Uslar cast steel. It is 4 feet long, nearly three inches wide and over half an inch thick, and weighs 25 pounds. It hangs from a 16-foot silk thread of a thousand strands\HistoricalNote{536}, fixed above the ceiling of the room and passing freely through a small opening made in the ceiling. The immediate reason for choosing so heavy a bar was to make it immune to drafts. These cannot always be kept entirely out of this room and have often interfered with the smaller apparatus, despite the box enclosing and protecting it. The result met \textit{this} expectation and far exceeded the others as regards the accuracy of all observations that can be made with it. Of course, the iron present in many parts of the astronomical observatory rules out absolute observations of declination and intensity with this instrument.

The greatest swing the box enclosing the bar allows is about 27 degrees. The greatest that can still be measured directly on the scale is 9 to 10 degrees, because with larger swings the telescope’s line of sight no longer meets the mirror, which is nearly four inches wide. Once the bar is set oscillating, the oscillations decrease in geometric progression so slowly that it often takes 10 hours or more for them to fall to half their size. Sometimes this happens much sooner, as we will discuss specifically below. The bar now suspended is the strongest of a larger number made for the physics cabinet. One oscillation takes a little over 42 seconds. This duration undergoes some very small changes owing to temperature and the variability of the Earth’s magnetism (and the strength of the bar itself may perhaps also change in time, although no change has yet been perceptible at all). A few oscillations suffice to determine the duration so precisely that one can leave the bar for 8 hours or more and afterwards be certain how many oscillations it has completed in the interval.

\Needspace{18\baselineskip}
\SourceLeaf{537}The electrodynamic experiments that can be made with this apparatus are just as interesting as the purely magnetic observations. For this purpose the bar is surrounded by a multiplier coil like that of the bar in the magnetic observatory, except that this coil is larger and has 2700 feet of wire in 270 turns. This multiplier coil has been connected to the large wire circuit mentioned in the earlier report, which connects the astronomical observatory, the magnetic observatory and the physics cabinet. In this circuit the galvanic current has to pass through a total wire length of 11000 feet, or nearly half a geographical mile, and then acts on three magnetic instruments at once, namely
\begin{list}{}{\setlength{\leftmargin}{2em}\setlength{\labelwidth}{2em}\setlength{\itemsep}{0pt}\setlength{\parsep}{0pt}}
\item[I.] the 25-pound bar in the astronomical observatory,
\item[II.] the 4-pound bar in the magnetic observatory\newline (multiplier coil of 200 turns)
\item[III.] the one-pound bar in the physics cabinet\newline (multiplier coil of 160 turns).
\end{list}
\noindent Individual parts of the circuit can easily be cut out at will in many different combinations.

Of the numerous experiments already made with these instruments, we give only a few here.

When a galvanic current is connected to the circuit, the magnetic bars in the three instruments begin to move so promptly that the start of their motion can be observed to within a small fraction of a second of time. Comparing the clocks at the three instruments gives results in such perfect agreement, whether the current is generated at one end, the other end or in the middle, that it fully confirms that the time the current takes to travel half a mile is too short to measure. Such a result might of course already have been anticipated from the interesting experiments of \textsc{Wheatstone}, recently published in the \textit{Philosophical Transactions} for 1834, according to which the electric current in metal seems to travel faster than light in space. Even so, those experiments did not warrant an unqualified inference about the behaviour of a \textit{galvanic} current and its effect on the magnetic needle.

The intensity of a galvanic current is measured or determined by the deflection of the magnetic needle, and thus in the first instance by scale divisions, though
\SourceLeaf{538}obviously in different units in the three instruments, depending on the dimensions of the multiplier coils and the value of the scale divisions in seconds of arc. Numerous experiments now show that the deflections produced in the three instruments by the same current at the same instant always stand in an exactly constant ratio, whether the current is generated at one end, at the other end or in the middle. This gives the important result that the current has the same intensity throughout its length, or at least loses none that can be detected. In future, particular attention will be paid to whether this result continues to hold under special circumstances, especially during heavy rain.

At all three instruments, commutators (gyrotropes) are connected to the circuit, making it easy to reverse the direction of the current. Court Councillor \textsc{Gauss} gave the commutator in the astronomical observatory a special design, so that a single press of the finger reverses the current quite instantaneously. Repeating this reversal at intervals always equal to the duration of one oscillation of one of the bars makes the oscillations of that bar grow larger and larger. This was used in an experiment that produces a striking mechanical effect. Professor \textsc{Weber} had a light release mechanism for an alarm clock or a pendulum clock fitted beside the magnetic bar in the physics cabinet. Each time, the current sent from the astronomical observatory operates the release perfectly after a couple of oscillations. It is obvious that the 25-pound bar could produce a much stronger mechanical effect.

These instruments are especially valuable in investigating the mathematical laws that govern the production and action of the magneto-electric induction discovered by \textsc{Faraday}, and in reducing it to absolute measures. Court Councillor \textsc{Gauss} will publish the results of his investigations on this subject elsewhere in due course. Of the devices used, we mention only one here, with which this induction can be produced in a way that is both simple and precisely measurable. A covered wire is wound around a wooden spool in 1050 turns, with its ends connected to the circuit through the commutator\SourceLeaf{539}. This spool can be slipped over the free half of a strong magnetic bar. Each time this is done, a galvanic current passes through the circuit, either strong and brief or weaker and longer-lasting, according as the operation is performed more quickly or more slowly, so that the total effect of sliding the spool on once is independent of the speed of the operation. The current itself always lasts only as long as the spool is moving. Pulling the spool off produces a current in the opposite direction, as does slipping it on with the opposite end first. If the movement is very quick, the effect of the current on the magnetic needle in one of the multiplier coils connected to the circuit can be treated as an instantaneous impulse of a definite strength. Pulling the spool off and putting it back on the other way round thus gives the magnetic needle two impulses in the same direction. Pulling it off again and then slipping it back on the other way round would give two impulses equal to each other and opposite to the preceding ones. If the commutator is reversed in between, however, the latter two effects are also in the same direction as the first two. One complete cycle of this kind (pulling the spool off, putting it back on the other way round and reversing the commutator) can easily be done in two seconds. During one oscillation of the large magnetic bar one can therefore, if desired, easily complete 21 cycles at a steady pace and set the bar moving so strongly that the whole scale passes out of the telescope’s field of view. This brief indication will suffice to show how the strength of the galvanic current produced by this kind of induction can be measured precisely. Its strength also depends, however, on the resistance offered by the circuit itself, and increases more or less as more or fewer sections of the circuit are cut out. In this way the ratios of the resistances in the individual parts of the circuit and in the multiplier coils have been determined with great precision. A wide variety of combinations have also given the fullest confirmation of the beautiful law stated by \textsc{Ohm}, which governs the intensity of a current when it divides. Nearly the same results have also been found with hydrogalvanic currents. These are less suited to such determinations because their strength varies, and in any event this still calls for special precautions in the experiments. It may be of some interest to note
\SourceLeaf{540}that the total resistance of the double wire connection carried through the air between the astronomical observatory and the physics cabinet, in a wire length of more than 6000 feet, is only about half the resistance the current encounters in the multiplier coil of the magnetic observatory alone (wire length 1100 feet), or only one sixth of the resistance in the whole circuit. The difference in wire thickness readily accounts for this, and all experiments confirm that for wires of the same metal the resistance is always directly proportional to the length and inversely proportional to the \textit{area} of the cross-section.

We mentioned above that the decrease in the arc of oscillation of the large needle was very unequal at different times. Similar differences had already appeared in 1832 with the small instruments, and later with the needle in the magnetic observatory. These differences always remain within far narrower limits, however, than in the case of the bar in the astronomical observatory. There the decrease in the arc from one oscillation to the next varied between $\frac{1}{5000}$ and $\frac{1}{50}$ in different series of experiments. This remarkable phenomenon drew the particular attention of Court Councillor \textsc{Gauss}. It seems to result from the concurrence of \textit{several} causes, some of which remain puzzling even now. In the meantime, Court Councillor \textsc{Gauss} has managed to identify the cause with by far the strongest influence. He observed that the arc of oscillation always decreased much faster when the circuit was closed than when it was open. It was then easy to recognise the cause of the rapid decrease as the reaction of a galvanic current produced in the circuit by induction through the oscillation of the needle itself. The current has the opposite direction in the next return swing and always acts to reduce the arc of oscillation. The explanation was fully confirmed. The arc decreased most slowly with the circuit open, more rapidly with the circuit closed but complete, still more rapidly when individual sections of the circuit were cut out, and most rapidly of all (so that the arc of oscillation fell to half in half an hour) when the circuit was closed immediately beyond the large bar’s multiplier coil. Indeed, these differences followed exactly the size of the part of the circuit left in operation.

\SourceLeaf{541}Once this explanation had been found, it was easy to foresee the outcome of some experiments that may well rank among the most striking in electromagnetism, and even to calculate the quantitative relations of the phenomena in advance. Repeated experiments confirmed them fully every time. They are as follows.

If the magnetic bar in the astronomical observatory (I) is set oscillating, with swings about as large as the box permits, these oscillations have no effect at all on the needles in the magnetic observatory (II) or in the physics cabinet (III). They \textit{remain} at rest if they were previously at rest, provided the circuit is open or at least the multiplier coils surrounding these latter needles are cut out of it. But the moment the circuit is closed, or, for example, the multiplier coil of II is brought into the closed circuit, needle II immediately begins to oscillate along with I. If needle II was already oscillating, its oscillations take on the distinctive character of \textit{mixed} oscillations. One depends on the initial state and has the same period as this needle’s oscillations under the influence of the Earth’s magnetism alone ($20''$), while the other follows a period of $42''$ (like the large needle I), and its size is proportional to the arc of oscillation of I (about $\frac{1}{200}$ when the circuit is cut off beyond the multiplier coil of II). This agrees completely with the results of theory. So does the circumstance, always confirmed exactly, that the oscillations of I and the induced oscillations of II, although they have periods of equal duration, do not begin together, but invariably differ in this respect by half an oscillation time ($21''$), and in the direction predicted by theory under the given circumstances. What has been said here by way of example about needle II occurs in a very similar way with needle III, whose natural oscillation time is $14''$, and which under the action of induction follows compound oscillations with periods of $14''$ and $42''$.

According to theory, the result must be quite different when a second needle with a multiplier coil is in the circuit in which the large bar oscillates, and the natural oscillation time of the second needle is exactly the same as that of the large magnetic bar. This second needle, which is quite still
\SourceLeaf{542}while the circuit is open, likewise begins to oscillate along with the bar the instant the circuit is closed. Its oscillations, of the same duration as the natural ones, increase steadily in size until (only after a very long time) they reach a maximum at which the resistance of the air balances the increase due to the force of induction. To carry out this remarkable experiment in practice (since it was impracticable at the time to suspend a large bar for lack of a second multiplier coil to suit it), the one-pound bar in the physics cabinet was made astatic in the usual way by coupling it with a similar, slightly less strongly magnetised bar. More precisely, it was made into a double needle whose natural oscillation time was brought to exactly $42''3$. The experiment succeeded perfectly with this arrangement. The bar oscillating in the astronomical observatory transmitted its oscillations to the double needle in the physics cabinet the moment the circuit was closed, as if by a wonderful sympathy, so that each successive oscillation was about 50 scale divisions or half a degree larger than the preceding one. Soon the entire image of the scale passed out of the field. Yet the increasing speed at which it passed through the field of view still made it possible to see the increase in the arc of oscillation. This wonderful play of sympathetic motion was observed for over an hour.

\Needspace{8\baselineskip}
It scarcely needs to be said that the four-pound bar in the magnetic observatory also induces a current in the closed circuit. Its existence is unmistakable from the rapid decrease in the arc of oscillation. It must therefore also act on the other two bars in a way similar to the effects shown by the first experiment. Calculation shows, however, and experience confirms, that these effects are too small to be perceptible. The weakest of the three bars would therefore be still less able to produce perceptible effects of this kind.

\clearpage
